Betrachten Sie die Funktion $f(z)$ mit
\[
f(x+iy)
=
e^x (\cos y + i \sin y)
=
\underbrace{e^x\cos y}_{\clap{$\displaystyle =u(x,y)$}}
\mathstrut + i
\underbrace{e^x\sin y}_{\clap{$\displaystyle =v(x,y)$}},
\]
\begin{teilaufgaben}
\item
Prüfen Sie nach, dass $u(x,y)$ und $v(x,y)$ die
Cauchy-Riemann-Differentialglei\-chun\-gen erfüllen.
\item
Zeigen Sie, dass die Funktion für reelle Argumente mit der
(reellen) Exponentialfunktion übereinstimmt und damit eine
komplexe Verallgemeinerung der Exponentialfunktion ist.
\item
Berechnen Sie die komplexe Ableitung aus $u$ und $v$ einerseits als
partielle Ableitung nach $x$ und andererseits als $1/i$ mal die
partielle Ableitung nach $y$ und überzeugen Sie sich, dass Sie
dasselbe erhalten.
\end{teilaufgaben}

\begin{loesung}
\begin{teilaufgaben}
\item
Die Ableitungen sind
\begin{equation}
\left.
\begin{aligned}
\frac{\partial u}{\partial x}
&=
\phantom{-}
e^x \cos y,
&
\frac{\partial u}{\partial y}
&=
-
e^x \sin y
\\
\frac{\partial v}{\partial x}
&=
\phantom{-}
e^x\sin y,
&
\frac{\partial v}{\partial y}
&=
\phantom{-}
e^x \cos y
\end{aligned}
\right\}
\Rightarrow
\left\{
\begin{aligned}
\frac{\partial u}{\partial x}&=\phantom{-}\frac{\partial v}{\partial y}
\\
\frac{\partial u}{\partial y}&=-\frac{\partial v}{\partial x}.
\end{aligned}
\right.
\end{equation}
\item
Für $y=0$
folgt
\[
f(x)
=
e^x(\cos 0 + i \sin 0)
=
e^x,
\]
wie behauptet.
\item
Die Ableitungen sind einerseits
\begin{align*}
\frac{df(z)}{dz}
&=
\frac{\partial u}{\partial x}+i\frac{\partial v}{\partial x}
\\
&=
e^x\cos y + ie^x \sin y
=
e^x(\cos y+i\sin y)
=
f(z),
\intertext{und andererseits}
&=
\frac{1}{i}
\frac{\partial u}{\partial y}
+\frac{\partial v}{\partial y}
\\
&=
-i
(-e^x\sin y) + e^x \cos y
=
e^x(\cos y+i\sin y)
=
f(z).
\end{align*}
Wie erwartet ist $f'(z)=f(z)$ und damit erfüllt $f(z)$ die
Differentialgleichung der Exponentialfunktion.
\qedhere
\end{teilaufgaben}
\end{loesung}


